\section{A weaker canonical target: growth rather than convergence (NS-66)}
\label{sec:ns66-canonical-growth}

Strong convergence of the full canonical sequence is excluded, but
that does not exclude slow growth or bounded subsequences of its norms.
One Mellin integration makes a complete Abel-summation estimate
available. The resulting growth exponent is an exact reformulation
of the zero-location question, not an independently established RH bound.

For a fixed integer $r\geq1$ define
\[
 D_{N,r}=\norm{J^r(\chi+B_N)}_H,\qquad
 \beta_* =\sup\{\Re z:\zeta(z)=0,\ 0<\Re z<1\}.
\]
Known critical-line zeros give $1/2\leq\beta_*\leq1$.

\begin{proposition}[Complete growth bounds for the smoothed canonical error]
\label{prop:ns66-growth-bounds}
If $z=\beta+i\gamma_*$ is a zero with $1/2<\beta<1$, then
\begin{equation}
 D_{N,r}\geq\frac{\sqrt{2\beta-1}}{|z|^{r+1}}N^{\beta-1/2}
 \quad\text{for every }N\geq1.
 \label{eq:ns66-zero-growth}
\end{equation}
Conversely, suppose for some $1/2\leq\theta<1$ and constant $C$
that $|M(u)|\leq Cu^\theta$ for all $u\geq1$. Then
\begin{equation}
 D_{N,1}\leq\sqrt2+
  \frac{C\norm{\sigma_1}}{1-\theta}N^{\theta-1/2}
  +C\sqrt\kappa\int_1^N u^{\theta-3/2}\,du,
 \qquad D_{N,r}\leq2^{r-1}D_{N,1}.
 \label{eq:ns66-mertens-upper}
\end{equation}
The integral is $(N^{\theta-1/2}-1)/(\theta-1/2)$ for
$\theta>1/2$, and $\log N$ for $\theta=1/2$.
The Mertens bound in this implication is a hypothesis, not a result
proved here.
\end{proposition}
\begin{proof}
Exact exterior interpolation gives $\chi+B_N=0$ on $x>1/N$.
Since $Jf(x)$ only uses arguments at least $x$, $J^r(\chi+B_N)$
is supported in $(0,1/N]$. Its Mellin value at $z$ is
$z^{-(r+1)}$: the zeta factor annihilates every basis contribution.
Cauchy--Schwarz on that support gives
\[
 |z|^{-(r+1)}\leq D_{N,r}
 \left(\int_0^{1/N}x^{2\beta-2}\,dx\right)^{1/2}
 =D_{N,r}\frac{N^{1/2-\beta}}{\sqrt{2\beta-1}},
\]
proving~\eqref{eq:ns66-zero-growth}.

Under the stated Mertens hypothesis, partial summation makes
$\sum\mu(n)/n$ convergent and identifies its value as
$\lim_{s\downarrow1}1/\zeta(s)=0$. Hence
\[
 \eta_N=-\int_N^\infty M(u)u^{-2}\,du,\qquad
 |\eta_N|\leq\frac{C}{1-\theta}N^{\theta-1}.
\]
The strong derivative $\partial_u\sigma_u=-\rho_u/u$ gives the
complete Bochner summation-by-parts identity
\begin{equation}
 JB_N=-N\eta_N\sigma_N+
           \int_1^N\frac{M(u)}u\rho_u\,du.
 \label{eq:ns66-bochner-abel}
\end{equation}
There is no lower-endpoint contribution, since $M(1^-)=0$.
Use $\norm{\sigma_N}=N^{-1/2}\norm{\sigma_1}$,
$\norm{\rho_u}=\sqrt{\kappa/u}$ and $\norm{J\chi}=\sqrt2$.
The triangle inequality yields~\eqref{eq:ns66-mertens-upper},
including all coefficient and tail-correction terms. Boundedness
of $J^{r-1}$ gives the final assertion.
\end{proof}

\begin{corollary}[The canonical norm-growth exponent records the rightmost zeros]
\label{cor:ns66-growth-exponent}
For every fixed integer $r\geq1$ the following limit exists and equals
\begin{equation}
 \lim_{N\to\infty}\frac{\log(1+D_{N,r})}{\log N}
                  =\beta_*-\frac12.
 \label{eq:ns66-exponent}
\end{equation}
In particular, RH is equivalent to the subpolynomial bound
\[
 \text{for every }\varepsilon>0,\qquad
 D_{N,r}=O_{r,\varepsilon}(N^\varepsilon).
\]
Even a bounded cofinal subsequence of $D_{N,r}$ would suffice to
prove RH. No such bounded subsequence or unconditional subpolynomial
bound is supplied, and boundedness is not claimed necessary under RH.
\end{corollary}
\begin{proof}
For each zero to the right of the line,
\eqref{eq:ns66-zero-growth} bounds the lower limit by
$\beta-1/2$. Taking the supremum gives $\beta_*-1/2$; when
$\beta_*=1/2$, the same lower limit is at least zero trivially.

If $\beta_*<1$, choose $\beta_*<\alpha<\theta<1$.
The classical zero-free-half-plane implication, stated precisely
in~\cite[Lemma 2.1]{BaezDuarteStrong}, gives convergence of
$\sum\mu(n)n^{-\theta}$. Partial summation of its bounded partial
sums gives $M(x)=O_\theta(x^\theta)$.
Proposition~\ref{prop:ns66-growth-bounds} then bounds the upper
limit by $\theta-1/2$. Let $\theta\downarrow\beta_*$.
This includes RH by choosing $\theta>1/2$ arbitrarily close to the
line, without claiming a bound at the endpoint $\theta=1/2$.

If $\beta_*=1$, the direct bounds $|s_N|\leq1+\log N$ and
$\norm{\sigma_n}=\norm{\sigma_1}/\sqrt n$ give
$D_{N,1}=O(\sqrt N\log(2N))$, so the upper limit is at most
$1/2$. This matches the preceding lower bound.
Thus~\eqref{eq:ns66-exponent} holds in every case.
The functional equation identifies $\beta_*=1/2$ with RH.
Finally~\eqref{eq:ns66-zero-growth} diverges along every cofinal
subsequence if any zero lies to the right of the line, proving the
bounded-subsequence implication.
\end{proof}

\paragraph{What the weaker target changes.}
For this explicit coefficient rule, norm convergence is false, while
subpolynomial norm growth is equivalent to RH. These are compatible
statements. The exact growth reformulation does not itself improve
our knowledge of $\beta_*$ or supply the missing Mertens estimate.
It gives a second precise research target: a direct subpolynomial
bound for the complete smoothed canonical norm, or a bounded cofinal
subsequence. Increasing the number of Mellin integrations changes
constants in~\eqref{eq:ns66-zero-growth}; the corollary is stated
only for one fixed order. No inference from a changing normalization
or finite table is used.
